\Qhead{1.8}{Устойчивость стационарных режимов автономных и неавтономных колебательных систем}{Устойчивость стационарных режимов автономных и неавтономных колебательных систем. Временные и спектральные методы оценки устойчивости.}
\Subfile{1.8.1}{sections/01_oscillations/q08/s01}
\Subfile{1.8.2}{sections/01_oscillations/q08/s02}

\recap{Устойчивость означает ограниченность малых отклонений, асимптотическая --- их затухание. Во времени изучают эволюцию возмущений или энергию; спектрально --- показатели и мультипликаторы. Автономный цикл сохраняет нейтральный сдвиг фазы.}
\board{Возврат амплитуды по наклону энергетического баланса; левая полуплоскость; единичный круг; критическая точка Найквиста.}
\oral{Почему спектр наблюдаемого сигнала не доказывает устойчивость породившего его режима?}
\bait{Можно оставить единичный мультипликатор автономного цикла: сдвиг времени переводит на ту же орбиту.}
\sources{Пятаков~\cite{pyatakov}, PDF-с.~51--52: спектральный анализ укороченной системы; с.~67--68: устойчивость и аттракторы. Энергетический критерий связан с~\Qref{1.2.2}; Dahleh и соавт.~\cite{dahleh}, §19.4 --- частотный критерий. Флоке и приведённый критерий Гурвица изложены как математическое дополнение, не как дословная запись курса.}
